\section{The same original Hom-section retains genuine isomorphisms}
\label{sec:scheme-invertible-hom-section-isomorphism}

\begin{theorem}[Actual whole Hom-section and geometric pullback isomorphisms]
\label{sc0-hom-section-iso:whole-map}
\lean{EnvelopeProperty.SchemeInvertibleHomSection.isIso_equiv_of_isIso,EnvelopeProperty.SchemeInvertibleHomSection.isIso_pullback_equiv_of_isIso_pullback}
\leanok
For any original invertible module $A$, any original module $B$ and
any genuine isomorphism $\varphi:A\to B$, its SAME actual whole
Hom-section is a genuine isomorphism from the original structure
module to the SAME tensor-dual module. More generally, if the SAME
original map becomes an isomorphism after ANY actual geometric
pullback, the SAME entire original Hom-section becomes an
isomorphism after that pullback. No frames, desired inverses or
unit/tensor/dual comparison certificates are supplied.
\end{theorem}
\begin{proof}
The actual Hom-section correspondence is the genuine adjunction of
tensoring by the SAME original invertible line, with actual unit
and tensor-object comparisons. Its value is the equivalence unit
followed by the inverse functor applied to the actual original map
and the genuine original left unitor. These are isomorphisms for
an actual original isomorphism, hence the SAME whole section is
an isomorphism.

For an actual pullback, first apply this assertion to the SAME
pulled original map. The proved entire original
Hom-section pullback identity relates that SAME section to the
pullback of the original section through actual native monoidal
unit/tensor comparisons and the genuinely derived original
right-dual evaluation comparison. Every comparison is a true
isomorphism. Cancel these actual isomorphisms from the ENTIRE
identity, and cancel the pullbacks of the genuine original unit
and tensor-object isomorphisms. The SAME pulled original section
is therefore an isomorphism, without assuming any of those
comparison equations or a frame coefficient.
\end{proof}

This converts the actual valuative canonical differential
isomorphism into unitness of its SAME original relative canonical
section, before comparing original rational coefficients. The
complete module pushforward and vanishing remain downstream.
